\section{Discussion} \label{sec:discussion}

This section synthesizes our empirical results, analyzing the scientific, survey, operational, information-theoretic, and physical implications of the TARS v1 framework.

\subsection{Scientific Implications: Isolating Epistemic Recovery Ambiguity} \label{sec:disc_epistemic}
Exoplanet vetting conflates two distinct sources of uncertainty: aleatoric (is the source physical—a planet or binary?) and epistemic (can the period recovery pipeline reliably trace the signal?). Traditional classifiers ignore this distinction, training on mixed features that optimize ranking without explaining why candidates fail. 

The RAI isolates epistemic uncertainty via graph topology, providing an interpretable, physics-grounded metric that decouples signal recovery stability from source-degeneracy classification. Rather than trying to resolve physical source degeneracies (such as background blending or eclipsing binaries) using light curve shapes alone, the RAI answers a different question: \textit{How reliably can we trace a unique transit ephemeris given the stellar activity profile and observation windows?} Isolating this uncertainty provides a clear, physical metric that directly scales with the falsifiability of the signal, preventing downstream classifiers from being biased by training catalog selection effects.

\subsection{Survey Implications: TESS and PLATO Integration} \label{sec:disc_survey}
Wide-field space missions generate vast quantities of time-series photometry, yielding tens of thousands of Threshold Crossing Events (TCEs) that exceed basic detection thresholds.
\begin{itemize}
\item \textbf{TESS Pipeline Application}: The TESS pipeline (e.g., SPOC) currently utilizes heuristic-based Robovetter-like rules or deep neural networks to vet candidates. Incorporating the RAI as a lightweight metadata field would allow the pipeline to flag active host stars or window-crossing aliases before full morphology vetting. This would significantly reduce the manual vetting workload.
\item \textbf{PLATO Pipeline Alignment}: The upcoming PLATO mission, with its long baseline observations (up to several years per field) and multiple camera groups, will observe stars under varying window configurations. High-cadence monitoring will yield complex periodograms. TARS's graph-theoretic model could scale dynamically to these longer baselines, utilizing the alias graph to prune false period branches and potentially identify true planets early. Future validation on simulated PLATO light curves will be required to confirm this capability.
\end{itemize}

\subsection{Operational Implications: Follow-Up Optimization} \label{sec:disc_operational}
High-resolution spectroscopic radial velocity (RV) follow-up and adaptive optics (AO) imaging are the primary bottlenecks in exoplanet confirmation. These resources are extremely limited.
\begin{itemize}
\item \textbf{Pre-Filtering and Ranking}: Currently, targets are prioritized based on simple signal-to-noise ratio (SNR) thresholds. However, a high SNR candidate on an active star can be a false periodogram lock (such as a rotation alias). The RAI provides a calibrated probability that measures the uniqueness of the period. By ranking candidates using $P(Y=1 \mid z_{\text{RAI}})$, observers can prioritize targets that have a unique, stable ephemeris.
\item \textbf{Integrating with Vespa and Triceratops}: Running complex Bayesian FPP models like Vespa or Triceratops requires significant computing hours per candidate. The computational cost suggests that the RAI could serve as a lightweight pre-screening metric before more computationally intensive validation methods such as Vespa or Triceratops. Candidates with high recovery ambiguity ($\text{RAI} \gg 1$) can be prioritized lower, saving CPU hours and focusing computational validation on stable, high-confidence candidates.
\end{itemize}

\subsection{Information-Theoretic Implications: Parsimony vs. Overfitting} \label{sec:disc_parsimony}
Our Conditional Mutual Information (CMI) audits and Forward Feature Selection sweeps (Section~\ref{sec:results_parsimony}) reveal a critical statistical trend: adding non-ambiguity features beyond a small core systematically degrades the Blind Validation Partition AUROC from $0.6943$ down to $0.5868$. Traditional vetting pipelines include up to dozens of correlated features that act as collinear noise on the independent Blind Validation Partition, degrading generalization. By restricting the classification model to a single, Z-score standardized index (the RAI), TARS v1 achieves stable, comparable performance to complex ensembles, demonstrating that transit vetting models can be substantially simplified while improving calibration.

\subsection{Physical Implications of Graph Entropy} \label{sec:disc_entropy}
Graph entropy ($x_{\text{ent}}$) measures the structural complexity and dispersion of the alias network. In quiet stars with stable transits, the period search yields a single node, resulting in a graph entropy of zero. On active stars, starspot groups cross the stellar disk, creating quasi-periodic dips. Because starspots grow, decay, and migrate across latitudes on timescales of days to weeks, the periodic modulation shifts phase and amplitude. This causes the transit search to lock onto multiple harmonics of the rotation period ($P_{\text{rot}}/2, 2 P_{\text{rot}}$, etc.) or window function beats. These harmonics form a highly connected, structured subgraph in the alias network. Graph entropy directly measures this physical complexity: a dense, highly structured subgraph yields a lower entropy than a scattered, unstructured set of candidates, linking the mathematical properties of the network directly to physical stellar photosphere activity.

\subsection{Comparative Analysis with Prior Classifiers} \label{sec:disc_comparison}
To highlight the parsimony and simplicity of TARS v1, we compare its structural characteristics with three primary baseline vetting pipelines:
\begin{itemize}
\item \textbf{Robovetter (Heuristic)}: Relies on hard, non-differentiable thresholds sensitive to calibration shifts in detrending. TARS v1, in contrast, translates continuous Z-score standardized features into a smooth logistic probability link for calibrated probabilistic grading of candidates near thresholds.
\item \textbf{Autovetter (Random Forest Ensemble)}: Autovetter trains a high-capacity random forest on 16 or more morphology and consistency features, where collinearity leads to decision tree overfitting and uncalibrated probabilities. TARS v1 eliminates these correlations by constructing a single parsimonious index.
\item \textbf{Astronet (Deep CNN)}: Astronet trains a deep convolutional neural network directly on 1D phase-folded light curves. While Astronet achieves high validation AUROC, it is a black-box model requiring GPU acceleration. TARS v1 requires no deep model training, runs in $< 0.1$ seconds on a single CPU core, and remains fully interpretable.
\end{itemize}

\subsection{Stewardship and Scientific Value of the TARS-250K-R1 Corpus} \label{sec:disc_corpus_value}
Rather than treating the TARS-250K-R1 operational corpus simply as a validation of pipeline throughput, we position this 250,010-target dataset as a foundational scientific infrastructure supporting several distinct avenues of future astronomical research:
\begin{enumerate}
\item \textbf{Stellar Ambiguity Atlas}: By computing the RAI for every recovered transit candidate, we compile a public ``ambiguity atlas'' mapping period, graph entropy, harmonic density, and recovery confidence. This provides the community with a reusable catalog to study how signal recovery degeneracies correlate with host star parameters (crowding, magnitude, temperature, and Galactic latitude).
\item \textbf{Targeted Follow-up Engine}: The pipeline filters the 250,010 stars down to 775 candidates, which are ranked by the RAI to yield the top 100 high-confidence discoveries. This acts as a targeted discovery engine, offering observers the highest-priority targets for manual vetting and ground-based spectroscopic radial velocity verification.
\item \textbf{Open Vetting Benchmarks}: Freezing these ambiguity scores establishes an open reference benchmark. Other transit detection search engines (e.g., box-fitting or deep learning models) can run on this standardized corpus to directly compare period-recovery rates, alias-rejection boundaries, and calibration reliability.
\item \textbf{Cross-Mission and Time-Series Analyses}: The RAI framework can be applied to Kepler, K2, and simulated PLATO datasets to study how ambiguity shifts across instruments. Furthermore, tracking the time-evolution of the RAI across multi-sector observations offers a diagnostic tool to study the stability of transit recovery under active starspot lifetime variations.
\end{enumerate}

\subsection{Engineering vs. Astrophysical Distinctions} \label{sec:disc_eng_vs_astro}
It is critical to distinguish the engineering accomplishments of this campaign from the resulting astrophysical claims:
\begin{itemize}
    \item \textbf{Engineering Conclusions}: The TARS pipeline can execute reproducibly at scale on a Ryzen 9 9950X architecture, achieving a sustained operational throughput of $9.13$ targets/sec with robust checkpointing, memory-bounded worker queue management, and automated fault isolation. No software crashes occurred during the production execution, and all data provenance anomalies were successfully captured under data availability rejections.
    \item \textbf{Astrophysical Claims}: The computed Recovery Ambiguity Index (RAI) is a mathematically consistent representation of alias graph complexity and period-recovery dispersion. However, while we have demonstrated that its five underlying features are not strongly collinear ($VIF < 2.50$) and correlate with manual vetting labels, we do not claim that the RAI has confirmed new exoplanets in this work. Resolving the physical nature of these candidates (e.g. distinguishing planets from astrophysical false positives such as grazing eclipsing binaries or blended stars) requires follow-up spectroscopy, radial velocity measurements, and catalog cross-matching.
\end{itemize}

\subsection{Threats to Validity} \label{sec:disc_threats_to_validity}

\subsubsection{Internal Validity}
*   \textbf{Data Availability Constraints}: Exclusions due to missing files ($1.75\%$) are derived directly from null-path records in the upstream database rather than execution crashes in the transit-vetting algorithm.
*   \textbf{FITS File Integrity}: The five isolated corrupt inputs represent a negligible data loss rate ($0.002\%$) and do not impact global statistical distributions.

\subsubsection{External Validity}
*   \textbf{Hardware and Storage Dependencies}: Operational throughput was measured under a specific Ryzen 9 9950X CPU, local SSD storage, and local SQLite database engine configuration. Performance characteristics may shift significantly under network-attached filesystems or distributed relational databases.
*   \textbf{Instrument Constraints}: The pipeline was evaluated exclusively on TESS Sector 1--14 light curves.

\subsubsection{Construct Validity}
*   \textbf{RAI Metric Formulation}: The RAI represents a specific, physical signed-sum configuration of five features. Alternative definitions of signal recovery ambiguity may exist and yield different relative rankings.

\subsubsection{Statistical Validity}
*   \textbf{Statistical Significance Sensitivity}: Due to the extremely large sample size ($N \approx 250\text{k}$), standard hypothesis tests (such as the Kolmogorov-Smirnov test) become hyper-sensitive, rejecting the null hypothesis under minor, practically negligible differences. Standardized effect sizes and Wasserstein distances are therefore reported alongside $p$-values to distinguish practical relevance from statistical detectability.

\subsection{Discussion of Cross-Matching Completeness \& Limitations (EXP-002)} \label{sec:disc_exp2}
The catalog cross-matching campaign executed in EXP-002 yielded an exoplanet recovery completeness of $R_{\text{planets}} = 0.44\%$ ($4 / 912$). While low, our diagnostic analysis provides a physical explanation for this result:
\begin{itemize}
    \item \textbf{Search Window Limitations}: More than a fifth ($21.69\%$) of the confirmed exoplanets near our targets have orbital periods exceeding the single-sector TESS baseline of $27.4\text{ days}$. Because TARS requires at least two transits to confirm a period, it has a physical limit of $\sim 13.7\text{ days}$ for robust period recovery.
    \item \textbf{Window-Function Systematics}: The remaining mismatches are consistent with expected window-function behavior in single-sector TESS observations. The central 1-day downlink gap in each sector creates a window beat at harmonics of $13.7\text{ days}$, which can dominate the BLS periodogram when the true signal is faint, causing the period search to lock onto systematic noise rather than the exoplanet's period. While the diagnostics support this interpretation, the present experiment does not uniquely identify the causal mechanism for every unmatched target.
    \item \textbf{Multi-Sector Stitching Motivation}: These results demonstrate that period recovery—not coordinate matching—is the dominant limitation of the pipeline. This strongly motivates future work (specifically EXP-003) to stitch multi-sector observations together, extending the baseline and removing downlink gap beats.
    \item \textbf{Nature of Novel Candidates}: The $240,613$ targets classified as \texttt{NOVEL\_CANDIDATE} represent an uncataloged candidate pool requiring additional astrophysical vetting, rather than confirmed discoveries. The vast majority of these targets are likely stellar variability, systematics, or period aliases of known objects that failed the matching gates.
\end{itemize}

\subsection{Discussion of Multi-Sector Transit Recovery (EXP-003)} \label{sec:disc_exp3}
The Experiment 3 results validate the central hypothesis that multi-sector stitching resolves period-recovery systematics:
\begin{itemize}
    \item \textbf{Resolution of Period Aliasing}: By combining observations across a multi-year baseline, the chronological stitcher successfully removes the downlink gap beat frequencies that contaminate single-sector periodograms. This is demonstrated by the $40.0\%$ absolute increase in exoplanet recovery completeness in the Cohort B pilot.
    \item \textbf{Limitations of Linear Ephemerides}: While multi-sector stitching increases the signal-to-noise ratio, it assumes a strictly linear transit ephemeris. Strong transit timing variations (TTVs) caused by planet-planet interactions over multi-year baselines can smear the transit signal across multiple phase bins, reducing BLS power. Future improvements will incorporate non-linear ephemeris models.
\end{itemize}

\subsection{Discussion of Injection--Recovery Sensitivity (EXP-004)} \label{sec:disc_exp4}
The Experiment 4 results quantify the detection efficiency of the frozen TARS pipeline under controlled signal injection conditions:
\begin{itemize}
    \item \textbf{Collinearity and Signal-to-Noise Ratio}: The multivariable logistic regression reveals that physical transit depth becomes non-significant ($p = 0.9999$) when SNR is included. This is a direct consequence of multicollinearity, as the recovered SNR is functionally proportional to transit depth:
    \begin{equation}
    \text{SNR} \propto \frac{\text{Depth}}{\sigma_N} \sqrt{N_{\text{in-transit}}}
    \end{equation}
    where $\sigma_N$ is the light curve noise level and $N_{\text{in-transit}}$ is the number of data points during transit. Thus, SNR acts as the primary physical descriptor of signal recoverability, absorbing the depth dependency. This interpretation is quantitatively supported by the Variance Inflation Factors (VIF), which show moderate collinearity between SNR ($\text{VIF} = 1.675$) and depth ($\text{VIF} = 1.658$), confirming that SNR shares substantial variance with depth.
    \item \textbf{Period-Dependent Completeness Limits}: The highly significant negative period coefficient ($\beta_2 = -0.0968, p < 10^{-15}$) indicates that longer periods reduce recovery probability, which physically corresponds to the fewer transits observed over a fixed sector duration, reducing the integrated SNR.
    \item \textbf{Limitations and Parameter Boundaries}: The measured strict completeness of $48.20\%$ represents an optimistic upper bound under simplified signal assumptions. The sensitivity limits apply strictly to the injected parameter space:
        \begin{itemize}
            \item Orbital periods between $0.5$ and $20.0\text{ days}$.
            \item Circular orbits without orbital eccentricity or precession.
            \item Trapezoidal transits assuming uniform stellar disks (no limb darkening).
            \item Solar-like stellar characteristics used for Keplerian orbital spacing.
        \end{itemize}
        In real observations, exoplanet searches are affected by transit timing variations (TTVs), transit duration variations (TDVs), eccentric orbits, and stellar activity (rotation, flares, spots), which degrade performance further.
\end{itemize}

\subsection{Discussion of Comparative Baseline Benchmark (EXP-005)} \label{sec:disc_exp5}

EXP-005 evaluated TARS Stages~2--3 as a standalone event-space detector against Vanilla BLS and TLS. The methodology (Section~\ref{sec:method_pipeline_flow}) defines the intended TARS deployment as a hybrid pipeline: a threshold detection stage (BLS or equivalent) identifies candidate transit events, and TARS performs event-space consistency checking and period refinement downstream. EXP-005 deviated from this by evaluating TARS without a threshold detection front-end. The results should be interpreted accordingly.

\subsubsection{What EXP-005 Establishes}

The experiment establishes a well-defined operating boundary. Above the single-event $3\sigma$ floor (transit depth $\gtrsim 2356$~ppm in the evaluated sample), TARS recovers injected periods with a median relative error of $0.0104\%$ at 2,879--9,603$\times$ lower computational cost than BLS and TLS. Below this floor, the event-space architecture cannot proceed because Stage~2 produces no TRANSIT\_LIKE events for Stage~3 to process. This is not a competitive deficit relative to BLS or TLS; it is an architectural constraint that was documented in the methodology prior to the experiment.

The 100\% false positive rate in the FP trials (10 non-injected light curves) provides additional evidence for the same conclusion: without an upstream threshold gate, the standalone Stage~2 detector generates TRANSIT\_LIKE events from pre-existing stellar variability, which Stage~3 interprets as candidates. This confirms that TARS requires an upstream detection stage and should not be deployed as an isolated search engine. That requirement is consistent with the preregistered hybrid design.

\subsubsection{Hybrid Architecture as the Intended System}

The preregistered EXP-005 design placed TARS downstream of BLS threshold detection. This document predates the experimental results and establishes that the hybrid pipeline was the intended deployment before any benchmark outcome was known. The standalone evaluation was an execution deviation from that design, not a retroactive reframing.

The appropriate future experiment---which EXP-005 as preregistered was designed to address---is: \textit{Does RAI vetting applied to BLS-detected candidates improve F1-score and ROC-AUC relative to unvetted BLS?} That comparison positions TARS correctly as an ambiguity resolution layer, not a replacement for threshold detection. We defer that experiment to future work.

\subsubsection{Implications for System Design}

These results suggest a complementary hybrid architecture:

\begin{enumerate}
    \item \textbf{Threshold detection stage} (BLS, TLS, or a future detector): produces transit candidates across the full depth range, including below the single-event floor, via coherent phase-folding.
    \item \textbf{Event-space vetting stage} (TARS Stages~2--3): applied to candidates that produce individually visible events above the noise floor; performs fast period confirmation, alias disambiguation, and RAI scoring.
\end{enumerate}

This division of labor exploits the complementary strengths of each approach: folding-based methods have higher sensitivity in the low-SNR regime; the event-space architecture is orders of magnitude faster and provides interpretable ambiguity scores in the high-SNR regime. Evaluating this hybrid system end-to-end is the natural next step following EXP-005.

\begin{figure*}[htbp]
\centering
\begin{tikzpicture}[node distance=2.2cm, auto, >=latex', thick,
    box/.style={rectangle, draw, text width=6.5cm, text centered, rounded corners, minimum height=4.5cm},
    line/.style={draw, ->, shorten >=1pt}]
    
    \node [box, fill=red!5] (conv) {
        \textbf{Conventional Vetting Evaluation} \\
        \vspace{0.2cm}
        \begin{itemize}
            \item Single-split training/testing
            \item Optimizes overall AUROC/F1
            \item Treats host stars as homogeneous
            \item Uncalibrated vetting thresholds
            \item High-risk of out-of-domain collapse
        \end{itemize}
    };
    
    \node [box, fill=green!5, right of=conv, node distance=8.5cm] (frame) {
        \textbf{Our Multi-Dimensional Protocol} \\
        \vspace{0.2cm}
        \begin{itemize}
            \item Cross-mission out-of-domain transfer
            \item Binned Expected Calibration Error (ECE)
            \item Cohort partitioning (active/quiet stars)
            \item Parameter-space noise stress sweeps
            \item Nested-regression redundancy checks
        \end{itemize}
    };
    
    \draw[<->, dashed, line width=1.5pt] (conv) -- node[above] {Comparison} (frame);
    
\end{tikzpicture}
\caption{Conceptual comparison between conventional single-split vetting evaluation and our multi-dimensional validation protocol. Standard evaluation masks domain shift, calibration decay, and cohort failure modes, whereas our protocol systematically exposes them.}
\label{fig:conventional_vs_framework}
\end{figure*}

\subsection{Methodological Comparison: Beyond AUROC in Vetting Evaluation} \label{sec:disc_beyond_auroc}

To demonstrate why standard single-metric evaluation protocols are insufficient for exoplanet Vetting systems, we present a comparison of standard evaluation practices against the multi-dimensional validation protocol implemented in this work (see Figure~\ref{fig:conventional_vs_framework} and Table~\ref{tab:evaluation_comparison}).

\begin{deluxetable*}{ll}[htbp]
\tablecaption{Methodological Comparison: Standard Vetting Evaluation vs. Multi-Dimensional Validation Protocol\label{tab:evaluation_comparison}}
\tablewidth{0pt}
\tablehead{
  \colhead{Standard Evaluation Practice} & \colhead{Insights Revealed by Multi-Dimensional Protocol}
}
\startdata
\textbf{Standard single-split AUROC} & \parbox{12.0cm}{Out-of-domain performance collapse. While standard splits show stable performance, cross-mission validation reveals a collapse on Kepler long-cadence data ($\text{AUROC} \approx 0.48$) due to cadence-induced transit-edge blurring.} \\
\hline
\textbf{Local test-set accuracy} & \parbox{12.0cm}{Hidden probability calibration decay. Single accuracy metrics mask the fact that predicted probabilities become severely miscalibrated under domain shifts (Kepler ECE $= 0.5212$ vs. TESS ECE $= 0.0966$).} \\
\hline
\textbf{Homogeneous target splits} & \parbox{12.0cm}{Physical stellar cohort limitations. Uniform splits conceal that convective granulation noise on evolved hosts ($\log g < 4.0$) deforms graph topologies, rendering the classifier no better than random on giant stars.} \\
\hline
\textbf{Single-metric benchmarking} & \parbox{12.0cm}{Robustness and operational boundaries. Evaluating on static datasets hides the threshold where high-frequency noise and missing cadences (stress-tested up to 30\% data loss) override signal ephemerides.} \\
\hline
\textbf{Independent feature listing} & \parbox{12.0cm}{Predictive redundancy under nested control. Standard feature importance listing misses collinearity and predictive redundancies, whereas our nested regression stress-testing (EXP-019) directly exposes that the RAI lacks independent predictive value once controlling for BLS SDE.} \\
\enddata
\end{deluxetable*}

This comparison demonstrates that a high AUROC on a local training/validation split does not guarantee scientific validity or operational utility. Operationalizing a multi-dimensional validation protocol—encompassing out-of-domain checks, nested controls, physical cohort limits, noise stress-testing, and calibration reliability analysis—is essential for exposing the hidden failure modes and structural redundancies of exoplanet vetting classifiers before they are deployed on production survey streams.


